\subsection{DET-13: an explicit positive transform and its trace-class family}

The preceding obstruction is specific to the signed scalar itself.  There is
an affirmative transformed construction, using a known Berezin kernel.  The
new comparison here is its exact expression in the source's $r=0$ determinant
and the resulting ordinary Fredholm family; positivity of the Berezin kernel
is an existing theorem \cite[Theorem 3.8, \S4.1 (4.1), \S4.4 (4.7)--(4.9)]{neretin}.

\begin{proposition}[DET-13]
Let $h(z,z')=z^*\operatorname{diag}(I_k,-1)z'$, $k\ge1$, and
$T(z)=-h(z,z)>0$.  For the actual $r=0$, $g=1$ signed Gram expression put
\[
 S(z,z')=|h(z,z')|^2-T(z)T(z').
\]
For every $a>0$ the function on negative projective lines
\begin{equation}
 F_a([z],[z'])=
 \left(1+\frac{S(z,z')}{T(z)T(z')}\right)^{-a}
 =\left(\frac{T(z)T(z')}{|h(z,z')|^2}\right)^a
 \label{eq:positive-gram-transform}
\end{equation}
is a continuous positive-definite kernel with diagonal one.  It is invariant
under $U(k,1)$ and is the Gram kernel of an explicit continuous family of unit
vectors $v_x$ in a separable Hilbert space $\mathcal H_a$.

For every finite positive Radon measure $\nu$ on this ball, the operator
\begin{equation}
 T_{a,\nu}=\int v_x\langle v_x,\,\cdot\,\rangle\,d\nu(x)
 \quad\hbox{satisfies}\quad
 T_{a,\nu}\ge0,\qquad T_{a,\nu}\in S_1,\qquad
 \operatorname{Tr}T_{a,\nu}=\nu(\mathbb B^k).
 \label{eq:positive-frame}
\end{equation}
It is a trace-norm limit of positive finite-rank sums on this same Hilbert
space.  Its ordinary determinant is
\begin{equation}
 \det(1-wT_{a,\nu})=
 \sum_{n=0}^{\infty}\frac{(-w)^n}{n!}
 \int_{(\mathbb B^k)^n}
 \det\big[F_a(x_i,x_j)\big]_{i,j=1}^n\,d\nu^n,
 \label{eq:positive-frame-det}
\end{equation}
with absolute locally uniform convergence.  For a finite atomic measure
$\nu=\sum_{j=1}^N b_j\delta_{x_j}$ with $b_j\ge0$, its determinant is exactly
\[
 \det\left[\delta_{ij}-w\sqrt{b_i b_j}F_a(x_i,x_j)\right]_{i,j=1}^N.
\]
No RH or zero input is used in this construction.
\end{proposition}

\begin{proof}
Every negative line has a unique representative $(x,1)$ with
$x\in\mathbb B^k=\{x:\|x\|<1\}$.  Scaling either original vector cancels
from the quotient in \eqref{eq:positive-gram-transform}, and preservation of
$h$ proves its simultaneous $U(k,1)$ invariance.  In these coordinates,
\[
 F_a(x,y)=
 \frac{(1-\|x\|^2)^a(1-\|y\|^2)^a}
 {|1-\langle x,y\rangle|^{2a}}.
\]
The denominator is nonzero.  To specify the Hilbert realization, let
$E=\ell^2(\mathbb N^k)$, $c_\alpha=(a)_{|\alpha|}/\alpha!>0$, where
$(a)_n=a(a+1)\cdots(a+n-1)$ and $(a)_0=1$, and set
\[
 u_x=(1-\|x\|^2)^{a/2}
       (\sqrt{c_\alpha}\,x^\alpha)_\alpha,
 \qquad v_x=u_x\otimes\overline{u_x}\in E\otimes\overline E.
\]
The binomial and multinomial expansions give
$\sum_\alpha c_\alpha\overline{x}^\alpha y^\alpha
=(1-\langle x,y\rangle)^{-a}$; they converge absolutely since
$\sum_j|x_jy_j|\le\|x\|\|y\|<1$.  Thus $\|u_x\|=\|v_x\|=1$ and
$\langle v_x,v_y\rangle=F_a(x,y)$.  This is the rank-one-ball version of
the cited Berezin construction, with the normalization needed here.
Take $\mathcal H_a$ to be the closed span of the $v_x$.  Continuity of the
kernel and its unit diagonal give norm continuity of $v_x$, since
$\|v_x-v_y\|^2=2-2F_a(x,y)$.  Kernel invariance defines a well-defined
unitary action $U_gv_x=v_{gx}$ on this span and then its completion: every
finite linear relation has the same squared norm after the assignment, and
$U_{g^{-1}}$ supplies the inverse.  In particular the transformed group
function is a genuine unitary matrix coefficient.

Write $P_x=v_x\langle v_x,\,\cdot\,\rangle$.  It is a positive rank-one
projection with trace norm one, and
$\|P_x-P_y\|_1\le2\|v_x-v_y\|$.  Hence its integral in $S_1$ exists for
finite $\nu$, is positive, and has trace $\int1\,d\nu$.  More concretely,
choose compact sets with remaining measure tending to zero.  On each compact
set, a finite Borel partition and representative points give sums
$\sum_j\nu(E_j)P_{x_j}$; their trace-norm error is bounded by the omitted
mass plus $\sum_j\nu(E_j)\sup_{x\in E_j}\|P_x-P_{x_j}\|_1$.  Uniform
continuity makes the latter tend to zero.  This both constructs the integral
and proves the asserted common-space approximation.  No norm compactness of
the whole ball is assumed.

For the determinant formula, first compress to the span of the first $m$
vectors of an orthonormal basis of $\mathcal H_a$.  The finite-dimensional
determinant expansion and the Leibniz formula under the integrals give its
$n$th exterior coefficient as
\[
 \frac1{n!}\int
 \|Q_mv_{x_1}\wedge\cdots\wedge Q_mv_{x_n}\|^2\,d\nu^n,
\]
where $Q_m$ is the orthogonal projection.  Indeed, for each coordinate minor,
expand both alternating determinants: Fubini factors the $n$ integrals, and
permuting the $n$ variables gives the factor $n!$.  Summing coordinate minors
is the finite Cauchy--Binet identity.  As $m\to\infty$, these integrands
converge to the full Gram determinant and are bounded by one, by the Gram
inequality.  The compressed operators converge in trace norm: this follows
from the same rank-one bound and dominated convergence.  Therefore DET-05
and dominated convergence yield \eqref{eq:positive-frame-det}.  The absolute
value of its $n$th term is at most
$|w|^n\nu(\mathbb B^k)^n/n!$, which justifies the entire-series limit.

For atomic $\nu$, define $V:\mathbb C^N\to\mathcal H_a$ by
$Ve_j=\sqrt{b_j}v_{x_j}$.  Then $T_{a,\nu}=VV^*$ and
$(V^*V)_{ij}=\sqrt{b_i b_j}F_a(x_i,x_j)$.  Restricting to their finite
ranges, or comparing their nonzero eigenvalues, proves the finite determinant
identity, including zero weights and repeated points.
\end{proof}

This gives an explicit family satisfying the positivity and approximation
parts of DET-05, after a nonlinear transformation of the actual signed
determinant.  It is not a claimed historical construction.  The choice of
finite measure and the Xi spectral identity remain substantive arithmetic
inputs.  In particular a putative equality with $G$ forces at least
$\nu(\mathbb B^k)=\operatorname{Tr}B$ and
$\iint F_a(x,y)^2\,d\nu(x)d\nu(y)=\operatorname{Tr}B^2$; all the higher
exterior coefficients must also agree.  These first two necessities follow
from the trace above and
$\operatorname{Tr}(P_xP_y)=|\langle v_x,v_y\rangle|^2$ by absolute Fubini.
Thus the bare signed kernel is excluded by DET-12, whereas transformed
positive families really exist.  The remaining Xi identification cannot be
replaced by either an impossibility assertion or positivity alone.

In Grothendieck's language this is the ordinary analytic trace/determinant
mechanism of cote~1/78--80 and its finite exterior-power approximants
(SGA~6, V, Definition~2.1, Example~2.2(b), Lemma~2.2.1).  The homogeneous
Hermitian domain and its unitary kernel are not, by this argument, an
arithmetic cohomology, an $\ell$-adic realization, or a crystal over
$\operatorname{Spec}\mathbb Z$.
